% page 508 of the facsimile (image p-507 of the scan)
% block 162.6 x 254.9 mm ; left margin 8.0 mm
\pgc{-12.700mm}{0.000mm}{
\l{}
\l{}
\l{}
\l{\cel{3}500.}
\l{}
\l{\sou{rismo}.\cel{2}Bloko de (en Britania) 480 od (en mondo, ecepte Bri-}
\l{\cel{3}tania) 500 folii de papero. - DEIS.}
\l{}
\l{\sou{risolar}. (\sou{koquarto}). Igar (kom exemplo : +kartofli - terpomi-}
\l{\cel{3}jus fritita) krakanta quale pano-krusto. - FI.}
\l{}
\l{\sou{rispektiva}. (\sou{aludante plura personi, plura kozi}). Qua koncernas}
\l{\cel{3}singla de li, relate a le cetera. - DEFIRS.}
\l{}
\l{\sou{ristornar}. (\sou{trans)}. I. (\sou{navig.}) Abrogar, nihiligar (interkon-}
\l{\cel{3}sente) asekur-akto navala konseque de cirkonstanci qui even-}
\l{\cel{3}tis ante la departo di la navo. - II. Diminutar la pekunio-}
\l{\cel{3}quanto ye qua onu asekuris la navo, kande la valoro di lua}
\l{\cel{3}omna kargaji esas fakte min granda kam expektita. - DFIRS.}
\l{}
\l{\sou{ritmo}. (\sou{muziko)}. Aparo inter-simetra di la tempi forta e di la}
\l{\cel{3}tempi febla, qua rieventas periodope en frazo muzikala, en}
\l{\cel{3}verso, en tamburo-frapado. - DEFIRS.}
\l{}
\l{\sou{ritornelo.}\cel{2}(\sou{muziko}).En la akompano di kanto, kurta frazo muzi-}
\l{\cel{3}kala ye qua onu igas preirata o sequata singla kupleto.-}
\l{\cel{3}DEFIS.}
\l{}
\l{\sou{rituo.}\cel{2}(\sou{liturgio)}\cel{2}La ceremonii di kulto. - DEFIRS.}
\l{}
\l{\sou{rivo}. Parto di la tero-sulo qua esas quaze la bordo di aquo-}
\l{\cel{3}extensajo, olua limito. - FS.}
\l{}
\l{\sou{riv}alo. La persono qua kontestas (ulo ad ulu) ; singla de la}
\l{\cel{3}personi qui inter-luktas por havar ulo, ulu. - DEFIS.}
\l{}
\l{\sou{rivero}. Aquo-fluo inter du rivi, qua fluas aden fluvio. -}
\l{\cel{3}EF.}
\l{}
\l{\sou{riveto}. (\sou{tekn.)}\cel{2}Klovo di qua la pinto esas abatebla ed aplas-}
\l{\cel{3}tebla per martelago (pos igir lu trairar plako) por fixi-}
\l{\cel{3}gar lu. - EF.}
\l{}
\l{\sou{rizo}. (\sou{bot.}) I. Planto cereala, graminea, kultivata en la}
\l{\cel{3}regioni di qui la atmosfero esas tre varma. - L.\cel{1}\sou{oryza}}
\l{\cel{3}\sou{sativa}. - II. La grani de ta planto. - DEFIRS.}
\l{}
\l{\sou{rizopodo}. (\sou{zool.}) Protozoo ek un celulo nuda, qua havas pro-}
\l{\cel{3}longuri (pseupodi) quin lu uzas kom organi lokomocala e}
\l{\cel{3}pren-ala. - DEFIS.}
\l{}
\l{\sou{robo.}\cel{2}Longa vesto qua decensas til la pedi e varias, pri lua}
\l{\cel{3}formo e lua naturo, segun la persono qua +weras lu, o}
\l{\cel{3}segun lua destineso. - DEFS.}
\l{}
\l{\sou{robineto}. (\sou{tekn.}) Ajuto muntita an la extremajo extera di}
\l{\cel{3}tubo di fonteno, di tanko, di cisterno, e trairata da kle-}
\l{\cel{3}fo qua, segun la sinso di lua rotacigeso, retenas o lasas es-}
\l{\cel{3}kapar la liquido. - FI.}
}
